\documentclass[11pt]{article}
\usepackage[margin=1in]{geometry}
\usepackage[T1]{fontenc}
\usepackage{lmodern}
\usepackage{microtype}
\usepackage{amsmath,amssymb,amsthm}
\usepackage{booktabs,array}
\usepackage{enumitem}
\usepackage[hidelinks]{hyperref}
\setlength{\emergencystretch}{2em}
\renewcommand{\topfraction}{0.95}\renewcommand{\textfraction}{0.05}\renewcommand{\floatpagefraction}{0.8}

% ---- notation
\newcommand{\R}{\mathbb{R}}
\renewcommand{\P}{\mathbb{P}}
\newcommand{\E}{\mathbb{E}}
\newcommand{\one}{\mathbf{1}}
\newcommand{\N}{\mathrm{N}}
\newcommand{\abs}[1]{\lvert #1\rvert}
\newcommand{\norm}[1]{\lVert #1\rVert}
\newcommand{\ind}[1]{\mathbf{1}\{#1\}}
\newcommand{\Cov}{\operatorname{Cov}}
\newcommand{\Var}{\operatorname{Var}}
\newcommand{\Corr}{\operatorname{Corr}}
\newcommand{\sd}{\operatorname{sd}}
\newcommand{\lmin}{\lambda_{\min}}
\newcommand{\yr}{y^{\mathrm r}}
\newcommand{\Xr}{X^{\mathrm r}}
\newcommand{\bstar}{\beta^*}
\newcommand{\tstar}{\theta^*}
\newcommand{\sigs}{\sigma_*}
\newcommand{\Sz}{S_0}
\newcommand{\sz}{s_0}
\newcommand{\Sig}{\Sigma_n}
\newcommand{\Om}{\Omega}
\newcommand{\pis}[1]{\pi(#1\mid y)}
\newcommand{\PIP}{\mathrm{PIP}}
\newcommand{\HPM}{\mathrm{HPM}}
\newcommand{\MPM}{\mathrm{MPM}}
\newcommand{\Shat}{\hat S}
\newcommand{\Zhat}{\hat Z}
\newcommand{\tSz}{\tilde S_0}
\newcommand{\rn}{\sqrt{\log n/n}}
\newcommand{\OP}{O_{\P}}
\newcommand{\oP}{o_{\P}}
\newcommand{\Dn}{\mathcal D_n}
\newcommand{\Hl}{\mathcal H}

% ---- statements
\theoremstyle{plain}
\newtheorem{thm}{Theorem}
\newtheorem{prop}[thm]{Proposition}
\theoremstyle{definition}
\newtheorem{asm}{Assumption}
\renewcommand{\theasm}{A\arabic{asm}}
\newtheorem{planned}{Planned result}
\renewcommand{\theplanned}{P\arabic{planned}}
\newcommand{\means}{\par\smallskip\noindent\textit{What it means.}\ }
\setlist[enumerate]{leftmargin=2em,itemsep=1pt,topsep=3pt}
\setlist[itemize]{leftmargin=1.5em,itemsep=1pt,topsep=3pt}

\begin{document}
\begin{center}
{\LARGE SMC-BVS: the model, the sampler,\\[3pt] the assumptions and the theorems}\\[8pt]
{\large A short version \quad$\cdot$\quad 24 November 2026}
\end{center}
\medskip

\noindent SMC-BVS is a method for Bayesian variable selection in the linear model when the number of predictors $p$ is
much larger than the sample size $n$. The model and the priors are standard. What is new is the sampler, which is built to
find the posterior mass when predictors are correlated, and a set of theorems about the posterior and about the model the
software reports. This note states each of these once: the model, the sampler, the assumptions, six proved results and
five planned ones. It contains no proofs.

% =====================================================================================================
\section{The model and the priors}

\paragraph{Data.} We observe a response $\yr\in\R^n$ and an $n\times p$ matrix of predictors $\Xr$; the superscript
means ``raw''. The method centres the response and every predictor and scales each to sample variance one. We write $y$
and $X$ for the standardised versions, so every column $x_j$ of $X$ has mean zero and $\norm{x_j}^2=n-1$, and
$y^\top y=n-1$.\footnote{A predictor that is constant in the sample is dropped first, and the prior then uses the number
of predictors kept. In Theorem~\ref{thm:post} a dropped predictor has inclusion probability $0$ and a model outside the
model space has posterior probability $0$, and $p$ in the assumptions counts all predictors.}

\paragraph{Model.} A model is a set $S\subseteq\{1,\dots,p\}$ of predictors, and $X_S$ holds its columns. Given $S$,
\[
  y\mid S,\theta_S,\sigma^2\ \sim\ \N_n(X_S\theta_S,\ \sigma^2I_n).
\]
There is no intercept because the data are centred. The model space $\Om$ is the set of all models with at most
$K=\min(p,\lfloor n/2\rfloor)$ predictors.

\paragraph{Priors.}
\[
  \pi(S)\ \propto\ \binom{2p-1}{\abs{S}}^{-1}\ (S\in\Om),\qquad
  \theta_S\mid S,\sigma^2\ \sim\ \N\big(0,\ n\sigma^2(X_S^\top X_S+I)^{-1}\big),\qquad
  \sigma^2\ \sim\ \mathrm{IG}(0.01,\,0.01).
\]
The model prior is the beta-binomial prior: each predictor enters independently with probability $q$, and
$q\sim\mathrm{Beta}(1,p)$; it is restricted to $\Om$. Going from $k$ to $k+1$ predictors costs a factor of about
$(k+1)/(2p)$ in prior odds. The coefficient prior is Zellner's $g$-prior with $g=n$, with a unit ridge added so that it
stays proper when columns are collinear. Nothing is left for the user to tune.

\paragraph{Posterior.} The coefficients and the variance integrate out in closed form. With $G_S=X_S^\top X_S$,
$z_S=X_S^\top y$ and $H_S=(n+1)G_S+I$, the marginal likelihood of $S$ is
\[
  m(S)\ \propto\ \Big(\frac{\det H_S}{\det(G_S+I)}\Big)^{-1/2}\big(0.02+y^\top y-n\,z_S^\top H_S^{-1}z_S\big)^{-(n/2+0.01)},
\]
and the posterior over models is
\[
  \pis{S}=\frac{\pi(S)\,m(S)}{Z},\qquad Z=\sum_{S\in\Om}\pi(S)\,m(S).
\]
We call $\pi(S)m(S)$ the \emph{score} of $S$ and $Z$ the \emph{evidence}. The usual summaries are the inclusion
probabilities $\PIP_j=\sum_{S\ni j}\pis{S}$, the highest posterior model $\HPM=\arg\max_S\pis{S}$, and the median
probability model $\MPM=\{j:\PIP_j>1/2\}$.

\paragraph{Truth.} The theorems assume that the data come from a sparse Gaussian linear model,
\begin{equation}\label{eq:model}
  \yr=\alpha^*\one_n+\Xr\bstar+\varepsilon,\qquad \varepsilon\sim\N_n(0,\sigs^2I_n).
\end{equation}
The true model is $\Sz=\{j:\beta^*_j\neq0\}$, of size $\sz$. Its members are the true predictors; the others are null
predictors. On the standardised scale the coefficients are $\theta^*_j=\beta^*_j\times(\text{sample s.d.\ of predictor }j)$.

% =====================================================================================================
\section{The sampler}

\paragraph{The idea.} The posterior is known up to the constant $Z$. The difficulty is to find where its mass sits in an
enormous model space. Correlated predictors make this hard in two ways. A true predictor and a close proxy give two
models with nearly the same score. And a null predictor that is correlated with a true one looks useful on its own,
because it carries part of the true predictor's signal; we say it \emph{leaks}.

Sequential Monte Carlo starts many models (\emph{particles}) from an easy distribution close to the posterior and moves
them to the posterior in small steps. Our easy distribution, the \emph{start law}, treats small blocks of correlated
predictors exactly. That settles the choice between a true predictor and its proxies inside a block. What the start law
cannot see is leakage between blocks. A greedy search, run at the end from the mode of the start law, removes leaked
predictors one at a time. This search uses no random numbers.

\paragraph{The steps.}
\begin{enumerate}
  \item \emph{Blocks.} Cluster the predictors by their sample correlations into blocks $C_1,\dots,C_J$ of at most
    $B=\min(\lfloor\sqrt n\rfloor,\lfloor\sqrt p\rfloor,16)$ predictors. Inside each block, compute
    $h_k(U)=m(U)/m(\emptyset)$ for every subset $U\subseteq C_k$ (at most $2^{16}$ subsets).
  \item \emph{Start law.} Set $\pi_0(S)\propto\pi(S)\prod_{k}h_k(S\cap C_k)$. This is the posterior one would have if the
    marginal likelihood factorised over blocks. A recursion over the blocks gives its normalising constant $\bar Z_0$,
    exact draws from it, and its mode $\tSz$.
  \item \emph{Path to the posterior.} Let $\Delta(S)=\log\{m(S)/m(\emptyset)\}-\sum_k\log h_k(S\cap C_k)$, the part of the
    log marginal likelihood that comes from interactions between blocks. The path is
    $\pi_\xi\propto\pi_0\,e^{\xi\Delta}$ for $0\le\xi\le1$, and $\pi_1$ is the posterior. A pilot run chooses a grid
    $0=\xi_0<\xi_1<\dots<\xi_T=1$, aiming to keep $90\%$ effective sample size at each step. The grid is then frozen.
  \item \emph{Particles.} Four independent \emph{islands} each draw $N=\max(2500,n,p)$ particles from $\pi_0$ (four times
    as many if the pilot's diagnostics ask for it). At step $t$ an island reweights particle $S_v$ by
    $\omega^{(t)}_v=\exp\{(\xi_t-\xi_{t-1})\Delta(S_v)\}$, resamples the particles in proportion to these weights, and
    moves each particle by Markov moves that leave $\pi_{\xi_t}$ invariant: a Gibbs update of every predictor, a random
    swap, swaps with correlated neighbours, and one proposal drawn afresh from $\pi_0$. At $\xi=1$ a few more rounds of
    moves are run, and the last four are used for the estimates.
  \item \emph{Greedy searches.} A search moves from the current model to its best neighbour (add one predictor, delete
    one, or swap one) as long as the score strictly increases, up to a cap on the number of moves. Two searches are run:
    one from the best model the particles found, with cap $\lceil\log_2(n+p+1)\rceil$, and one from the start mode $\tSz$,
    with cap $\lceil\log_2(n+p+1)\rceil+\abs{\tSz}$.
  \item \emph{Report.} The \emph{reported model} $\Shat$ is the model with the highest score among all models the run
    scored. It is \emph{certified} if it is the end point of one of the two searches and no neighbour has a higher
    score: a local optimum under additions, deletions and swaps. The run also reports inclusion probabilities (averaged
    over particles and islands), the list of the $40p+1$ best models scored, and the evidence estimate
    \begin{equation}\label{eq:zhat}
      \Zhat=\frac14\sum_{r=1}^{4}\Zhat_r,\qquad
      \Zhat_r=m(\emptyset)\,\bar Z_0\prod_{t=1}^{T}\Big\{\frac1N\sum_{v=1}^N\omega^{(t)}_v\Big\}\quad\text{for island }r.
    \end{equation}
\end{enumerate}
When $\Om$ is small enough the method enumerates it instead, and everything it reports is exact (and $\Shat$ counts as
certified). A run can also stop without output, for instance when it would exceed its time or memory allowance. We say
that a run \emph{produces output} when it does not stop.

\paragraph{Very large $p$.} For $p$ near $10^6$ a parallel version runs $16$ islands as separate processes that do not
communicate, merges what they found at the end, and then runs the two greedy searches; it scores models with a faster
batched routine and re-scores every accepted model exactly. The theorems below are stated for the standard version, not
for this one.

% =====================================================================================================
\section{The assumptions}

Table~\ref{tab:assumptions} lists the assumptions, the theorems that use them, whether they are reasonable, and how they
compare with Song and Liang (2015). Each assumption is then stated once. Constants called \emph{fixed} do not depend
on $n$; everything else ($p$, the design, the coefficients) may change with $n$.


\begin{table}[!t]
\centering\footnotesize
\setlength{\tabcolsep}{3.5pt}
\renewcommand{\arraystretch}{1.08}
\caption{The assumptions at a glance. ``Used by'' gives theorem numbers. The last column compares with Song and Liang
(2015): \emph{same}, \emph{weaker} (ours asks less), \emph{stronger} (ours asks more), \emph{mixed}, or \emph{none} (no
analogue).}\label{tab:assumptions}
\begin{tabular}{@{}>{\raggedright\arraybackslash}p{0.215\linewidth}>{\raggedright\arraybackslash}p{0.075\linewidth}>{\raggedright\arraybackslash}p{0.375\linewidth}>{\raggedright\arraybackslash}p{0.275\linewidth}@{}}
\toprule
Assumption & Used by & Reasonable? & vs Song and Liang \\
\midrule
\multicolumn{4}{@{}l}{\textbf{The data model}} \\
A1 Gaussian linear model & 1, 4, 6(b) & the standard model; true by construction in our simulations & same (their model (1)) \\
A2 $p\le n^A$ & 1, 3, 6(b) & any polynomial $p$; $p=10^6$ at $n=200$ is $A\approx2.6$ & mixed: they allow $\log p$ up to a power of $n$; we allow bounded $p$ \\
A3 fixed true model size & 1, 6(b) & few strong effects; nearly necessary for the rate (Theorem~\ref{thm:sharp}(a)) & same as their strict sparsity (4) \\
\midrule
\multicolumn{4}{@{}l}{\textbf{The predictors and the signal}} \\
A4 law of the rows & 1, 3 & Gaussian or sub-Gaussian; bounded features are included, rare binary features are not & stronger than their Condition~1 \\
A5 fixed signals & 1, 3 & effects that do not shrink with $n$; Theorem~\ref{thm:dist} needs none & same, if their $\bstar$ is fixed \\
A6 eigenvalue floor & 1, 3 & holds in all our designs (floors $0.1$ to $1$); no upper bound, so common factors are allowed & mixed: stronger than their Condition~3 where the two can be compared \\
A7 weaker setting of Theorem~\ref{thm:dist} & 2 & A1--A4 relaxed, no lower bound on the effects; only exact duplicates and exact linear dependences fail & stronger in one respect: their Theorems 1--2 have no collinearity condition \\
\midrule
\multicolumn{4}{@{}l}{\textbf{The reported model (Theorem~\ref{thm:reported} only)}} \\
A8 visibility & 3 & holds with aligned signs or generic coefficients; can fail by cancellation & stronger than their Condition~5 (marginal, no margin) \\
A9(a) small leakage set & 3 & holds for block and autoregressive designs; fails for pervasive factors & none \\
A9(b),(c) floor and local regressions & 3 & follow from a floor on $\Sig$ and local correlation & floor as A6; locality none \\
\midrule
\multicolumn{4}{@{}l}{\textbf{Certified screening (Theorem~\ref{thm:screen} only)}} \\
A10 floor on sets of size $k_n+\sz$, signal bounds & 4 & fixed design; common factors allowed; the sample-size condition is conservative & as A6 and A5, on other sets \\
\midrule
\multicolumn{4}{@{}l}{\textbf{Not statistical}} \\
C1--C3 idealised computer & 3--5 & the usual idealisation in a theorem about an algorithm & their sampler statements are idealised limits too \\
\bottomrule
\end{tabular}
\end{table}

\subsection*{The data model}

\begin{asm}[Gaussian linear model]\label{a:gl}
The data follow~\eqref{eq:model}, with a fixed noise variance $\sigs^2>0$ and noise independent of the predictors.
\end{asm}

\begin{asm}[Dimension]\label{a:dim}
$2\le p\le n^A$ for a fixed $A\ge1$.
\end{asm}

\begin{asm}[Sparsity]\label{a:sp}
The number of true predictors $\sz\ge1$ is fixed.
\end{asm}

\subsection*{The predictors and the signal}

\begin{asm}[Law of the rows]\label{a:rows}
The rows of $\Xr$ are independent and identically distributed, with mean $m_n$ and covariance matrix $\Sig$, which is a
correlation matrix. In addition, one of:
\begin{enumerate}[label=(\alph*)]
  \item \emph{sub-Gaussian on sparse directions:} for a fixed $L>0$, a row $x$ satisfies
    $\E\exp\{\lambda u^\top(x-m_n)\}\le\exp\{L^2\lambda^2u^\top\Sig u/2\}$ for every $\lambda\in\R$ and every $u\in\R^p$
    with at most $5(A+1)(\sz+1)\log n$ nonzero entries;
  \item \emph{Gaussian:} the rows are $\N_p(m_n,\Sig)$. This is case (a) with $L=1$.
\end{enumerate}
\end{asm}

\begin{asm}[Fixed signals]\label{a:fs}
One of:
\begin{enumerate}[label=(\alph*)]
  \item \emph{raw scale:} $\alpha^*$ and $\bstar$ are deterministic, and $b_{\mathrm{lo}}\le\abs{\beta^*_j}\le
    b_{\mathrm{hi}}$ for every $j\in\Sz$, with fixed $0<b_{\mathrm{lo}}\le b_{\mathrm{hi}}<\infty$;
  \item \emph{standardised scale:} $\yr=\alpha\one_n+X\tstar+\varepsilon$, where $X$ is the standardised design,
    $\varepsilon\sim\N_n(0,\sigs^2I_n)$ is independent of $\Xr$, $\sigs^2>0$ is fixed, $\tstar$ is deterministic with a
    fixed number $\sz\ge1$ of nonzero entries, and $\theta_{\mathrm{lo}}\le\abs{\theta^*_j}\le\theta_{\mathrm{hi}}$ for
    every $j\in\Sz$, with fixed $0<\theta_{\mathrm{lo}}\le\theta_{\mathrm{hi}}<\infty$.
\end{enumerate}
\end{asm}
\noindent Form (b) says that the true coefficients are fixed numbers on the scale the method works on. It contains the
linear model and the sparsity assumption, so it takes the place of A1 and A3 in Theorem~\ref{thm:reported}. When the
predictors have unit population variance the two scales differ only by sampling error.

\begin{asm}[Eigenvalue floor]\label{a:efs}
$\lmin(\Sigma_{n,S})\ge\varphi^2$ for every set $S$ of at most $5(A+1)(\sz+1)\log n$ predictors, for a fixed
$\varphi\in(0,1]$. Here $\Sigma_{n,S}$ is the sub-matrix of $\Sig$ for the predictors in $S$.
\end{asm}
\noindent No small group of predictors is close to linearly dependent in the population. A floor on the whole matrix,
$\lmin(\Sig)\ge\varphi^2$, is enough. There is no upper bound on the eigenvalues.

\begin{asm}[The weaker setting of Theorem~\ref{thm:dist}]\label{a:dist}
The observations are $n$ independent copies of a pair $(x,y)$ with $x\in\R^p$. Fix an integer $s\ge0$ and constants
$A\ge1$, $L\ge1$, $c_\psi>0$, $\bar b>0$. For a level $M\ge1$, let $q_n=\min\{p,\lceil C_qM\log n\rceil\}$ with
$C_q=2(s(A+2)+4)/\log2$.
\begin{enumerate}[label=(\alph*)]
  \item \emph{Model.} $y=\alpha^*+x^\top\bstar+\varepsilon$ with $\varepsilon\sim\N(0,\sigs^2)$ independent of $x$ and
    $\sigs^2>0$. The quantities $\alpha^*$, $\bstar$, $\sigs$ and the law of $x$ may depend on $n$.
  \item \emph{Sparsity and dimension.} At most $s$ entries of $\bstar$ are nonzero, and $2\le p\le n^A$.
  \item \emph{Sub-Gaussian predictors.} Every predictor has positive variance. With
    $\zeta_j=(x_j-\E x_j)/\sd(x_j)$, $\E\exp(\lambda u^\top\zeta)\le\exp\{L^2\lambda^2\Var(u^\top\zeta)/2\}$ for every
    $\lambda\in\R$ and every $u$ with at most $q_n$ nonzero entries.
  \item \emph{No near-exact collinearity.} $\lmin(\Corr(x_T))\ge c_\psi/n$ for every set $T$ of at most $q_n$ predictors.
  \item \emph{Bounded signal.} $\abs{\beta^*_j}\,\sd(x_j)\le\bar b\,\sigs$ for every $j$.
\end{enumerate}
\end{asm}
\noindent Items (a)--(c) are A1--A4 with less asked: the noise level and the coefficients may change with $n$, and no
coefficient has to be large, or even nonzero. The floor in (d) tends to zero, so it excludes only predictors that are
almost exactly collinear.

\subsection*{The reported model}

The start law sees interactions only inside blocks of at most $16$ predictors. The next two assumptions are therefore
about how predictors relate once up to $15$ others are held fixed. Let $W\sim\N_p(0,\Sig)$ be a centred copy of one row
of the design. For a set $U$, $\Cov(W_j,W_l\mid W_U)$ is the partial covariance given the predictors in $U$, and
$\bar\beta_{U,l}=\Sigma_{n,UU}^{-1}\Sigma_{n,Ul}$ is the population regression coefficient of $W_l$ on $W_U$.

\begin{asm}[Visibility]\label{a:vis}
For every true predictor $j\in\Sz$ and every set $U\not\ni j$ of at most $15$ predictors,
$\abs{\Cov(W_j,\,W^\top\tstar\mid W_U)}\ge a$, for a fixed $a>0$.
\end{asm}
\noindent Each true predictor stays related to the signal $W^\top\tstar$, by a fixed margin, whichever $15$ other
predictors are held fixed. It fails when other predictors cancel a true predictor's signal: two true predictors with
coefficients $1$ and $-1$ and correlation $\rho$ each have covariance $1-\rho$ with the signal, which vanishes as
$\rho\to1$.

\begin{asm}[Small leakage set]\label{a:leak}
Fix $b\in(0,\tfrac12)$ and $b_H>0$ with $15b_H+b<\tfrac12$.
\begin{enumerate}[label=(\alph*)]
  \item \emph{Leakage set.} There is a deterministic set $Q\supseteq\Sz$ with $\abs{Q}\le n^b$ such that:
    \begin{enumerate}[label=(\roman*),leftmargin=1.8em]
      \item \emph{no leak outside $Q$:} $\abs{\Cov(W_l,W_j\mid W_U)}\le\varpi_n$ for every $l\notin Q$, every set $U$ of
        at most $15$ predictors and every $j\in\Sz\setminus U$, where $\varpi_n\sqrt{n/\log n}\to0$;
      \item \emph{locality:} every $l\notin Q$ has a deterministic neighbourhood $\Hl_l\supseteq\{l\}\cup\Sz$ with
        $\abs{\Hl_l}\le n^{b_H}$ such that, for every $U\not\ni l$ with $\abs{U}\le15$ and $U_0=U\cap\Hl_l$,
        $\abs{\Cov(W_l,W_v\mid W_{U_0})}\le\delta_n$ and $\abs{\Cov(W_j,W_v\mid W_{U_0})}\le\delta_n$ for every
        $v\in U\setminus\Hl_l$ and $j\in\Sz\setminus U_0$, where $\delta_n\to0$;
      \item \emph{room under the size cap:} for large $n$, $\abs{Q}\le K$, and either $Q$ is the set of all predictors or
        $K\ge2\abs{Q}+32\sz+1$.
    \end{enumerate}
  \item \emph{Floor.} $\lmin(\Sigma_{n,Q\cup\{l\}})\ge\varphi^2$ for every predictor $l$, with $\varphi$ as in A6.
  \item \emph{Local regressions.} There are sequences $d_n=o(\log n)$ and $L_n\to0$ with
    $L_n^2\log(\abs{Q}+1)=o(\log n)$, and for every predictor $l$ a deterministic set $N_l\subseteq Q\setminus\{l\}$ of at
    most $d_n$ predictors, such that $\norm{\bar\beta_{T,l}-\bar\beta_{T\cap N_l,l}}_1\le L_n$ for every $T$ with
    $\Sz\subseteq T\subseteq Q$ and $l\notin T$. (The second coefficient vector is filled with zeros to the length of the
    first.)
\end{enumerate}
\end{asm}
\noindent In words: (a) only a small set $Q$ of predictors, fewer than $\sqrt n$, can imitate part of the signal; every
predictor outside $Q$ is nearly uncorrelated with the true ones given any $15$ others, and conditioning on far-away
predictors changes nothing that matters. (b) No predictor is close to a linear combination of the members of $Q$. (c)
The regression of any predictor on a model between $\Sz$ and $Q$ is close to its regression on a few members of $Q$.
For independent predictors $Q=\Sz$. For block designs $Q$ is the union of the blocks that contain true predictors, and
for autoregressive designs it is a neighbourhood of size $C\log n$ around them. With a pervasive common factor every
null predictor is correlated with every true one, so $Q$ would have to be everything, and A9 fails once $p$ is larger
than $\sqrt n$.

\subsection*{Certified screening}

\begin{asm}[Conditions for certified screening]\label{a:scr}
The design is fixed. Let $\hat R_S=X_S^\top X_S/(n-1)$ be the sample correlation matrix of the predictors in $S$. For a
\emph{size budget} $k_n\in\{0,\dots,K-1\}$:
\begin{enumerate}[label=(\alph*)]
  \item $\lmin(\hat R_S)\ge\kappa$ for every $S$ with $\abs{S}\le k_n+\sz$;
  \item $\min_{j\in\Sz}\abs{\theta^*_j}\ge\underline\theta>0$ and $\norm{\tstar}_1\le M_1$;
  \item $(k_n+1)\log p+\log n\le c\,n$, with $c=\kappa^2\underline\theta^2/\{200(M_1^2+\sigs^2)\}$.
\end{enumerate}
The constants $\kappa>0$, $\underline\theta$, $M_1$ and $\sigs^2$ are fixed.
\end{asm}
\noindent These are an eigenvalue floor and signal bounds of the same kind as A6 and A5, now on the observed design and
on sets of the budget size. Item (c) lets the budget be as large as a small multiple of $n/\log p$. The true model size
$\sz$ need not be fixed.

\subsection*{The idealised computer}

Theorems~\ref{thm:reported}--\ref{prop:unbiased} are about runs of the software. They idealise the computer in three
ways. These are not assumptions about the data.
\begin{itemize}
  \item[\textbf{C1}] \emph{The method.} Exact arithmetic, the default settings, and a valid input and seed.
  \item[\textbf{C2}] \emph{The machine.} C1, and a machine that never overflows or fails, ideal random numbers, and at
    least the memory and time that the method's own formulas ask for. (Two hard-coded integer limits on the search caps
    are also ignored; no real input can reach them.)
  \item[\textbf{C3}] \emph{The sampler.} Exact arithmetic; ideal random numbers that, in the final run, are independent
    of the data and of the pilot; exact draws from the start law; and Markov moves that do what their definitions say.
\end{itemize}

% =====================================================================================================
\section{The theorems}

\subsection{Proved}

Theorems~\ref{thm:post}, \ref{thm:dist} and~\ref{thm:sharp} are about the exact posterior.
Theorems~\ref{thm:reported} and~\ref{thm:screen} and Proposition~\ref{prop:unbiased} are about what the software
returns.

\begin{thm}[Posterior concentration]\label{thm:post}
Assume A1, A2, A3, A4(a), A5(a) and A6. Then, under the joint law of the design and the noise,
\[
  1-\pis{\Sz}=\OP\big(\rn\big),\qquad
  \max_{1\le j\le p}\big|\PIP_j-\ind{j\in\Sz}\big|=\OP\big(\rn\big),
\]
and $\P(\HPM=\MPM=\Sz)\to1$.
\end{thm}
\means The posterior puts all but about $\rn$ of its mass on the true model, every inclusion probability is within that
distance of $0$ or $1$, and both standard point estimates are the true model. This holds for sub-Gaussian designs with
any correlation structure that has an eigenvalue floor, common factors included, and for any polynomial $p$.

\begin{thm}[Consistency of the fitted distribution]\label{thm:dist}
A posterior draw $(S,\beta,\sigma^2)$, in the original units, defines the conditional law
$\N(\bar y+(x-\bar x)^\top\beta,\sigma^2)$ of $y$ given $x$, where $\bar x$ and $\bar y$ are the sample means. Let $f$
be the joint density of $(x,y)$ when $x$ has its true law and $y$ given $x$ has this fitted law, let $f^*$ be the true
joint density, and let $d_H(f,f^*)=(1-\int\!\sqrt{ff^*})^{1/2}$ be the Hellinger distance.

There are constants $c_H>0$, $M_0$ and $n_0$, depending only on $(s,A,L,c_\psi,\bar b)$, such that the following holds.
Fix $M\ge M_0$ and assume A7 with this $M$. Then for every $n\ge n_0$ with $\sqrt n/\log n\ge M$,
\[
  \P^*\Big[\ \Pi\Big(d_H(f,f^*)>\sqrt{M\log n/n}\ \Big|\ \Dn\Big)>n^{-c_HM}\ \Big]\ \le\ n^{-c_HM},
\]
where $\Dn$ is the data set, $\P^*$ its law and $\Pi(\cdot\mid\Dn)$ the posterior of $(S,\beta,\sigma^2)$. If the
posterior is not defined, because fewer than two predictors are left or the response is constant in the sample, the
inner probability is read as $1$.
\end{thm}
\means The distribution that the posterior fits to $y$ given $x$ is within $\sqrt{M\log n/n}$ of the truth, except on
a set of posterior mass $n^{-c_HM}$ and a set of data sets of the same probability. The bound holds for each large $n$,
not only in the limit. It needs no lower bound on the coefficients: effects too weak to be selected do no harm to the
fitted distribution. Theorem~\ref{thm:post} is about selection; this one is about prediction.

\begin{thm}[The reported model equals the true model; simplified, full version in spec.tex, Theorem~3]\label{thm:reported}
Let $\Sig$ be positive definite and assume A2, A4(b), A5(b), A6, A8, A9 and C1. Then there are events $\mathcal G_n$,
determined by the data alone, with $\P(\mathcal G_n)\to1$, on which every run that produces output reports
$\Shat=\Sz$, whatever its seed and whatever its pilot chose. In particular
\[
  \P\big(\text{the run produces output and }\Shat\neq\Sz\big)\to0 .
\]
Under C2 every run produces output, for every $n\ge3$ that the method accepts, and so $\P(\Shat=\Sz)\to1$.
\end{thm}
\means Theorems~\ref{thm:post} and~\ref{thm:dist} describe a posterior that nobody can compute exactly at this size.
Theorem~\ref{thm:reported} is about the model the software prints: for almost every data set it is exactly the true
model, and this does not depend on the seed. The proof follows the design of the sampler. By visibility, the mode of
the start law contains every true predictor. Anything else in it is a leaked predictor, and so belongs to $Q$. The
greedy search from that mode deletes the leaked predictors one at a time and stops at $\Sz$, and no other model the run
scores beats $\Sz$. The search is deterministic, which is why the seed does not matter. The result covers independent,
block and autoregressive designs. It does not cover pervasive common factors, because A9 fails there.

\begin{thm}[Certified screening]\label{thm:screen}
Let the design be fixed and assume A1, A10 with size budget $k_n$, and C1. Then
\[
  \P\big(\text{the run produces output},\ \Shat\text{ is certified},\ \abs{\Shat}\le k_n,\ \text{and }\Sz\not\subseteq\Shat\big)\to0 .
\]
\end{thm}
\means A certified model is a local optimum of the posterior: no addition, deletion or swap of one predictor improves
it. The theorem says that a local optimum with at most $k_n$ predictors cannot miss a true predictor, except with
vanishing probability. If it missed one, bringing that predictor in would raise the score by more than chance can
explain. So a small certified model can be read as a screening set: it contains every true predictor. The theorem needs
neither visibility nor a leakage set, so it covers designs with strong common factors, where Theorem~\ref{thm:reported}
does not apply. It bounds the probability of one joint event. It does not say that a run will certify a small model, it
is not a statement conditional on certification, and a certified model may contain null predictors.

\begin{prop}[The evidence estimate is unbiased; simplified, full version in spec.tex, Proposition~5]\label{prop:unbiased}
Assume C3, and fix any data set on which the method uses the sampler and the pilot succeeds. Let $\mathcal C$ be
everything the pilot chose. Then the evidence estimate that each island would produce if run to the end satisfies
\[
  \E\big[\Zhat_r\mid\mathcal C\big]=Z\quad(r=1,\dots,4),\qquad\text{and hence}\qquad \E\big[\Zhat\mid\mathcal C\big]=Z ,
\]
where the expectation is over the random numbers of the final run. No condition on the design is needed.
\end{prop}
\means The grid and the number of particles are tuned on the data by the pilot. Because they are frozen before the
final run draws a single random number, the classical unbiasedness of the sequential Monte Carlo estimate of a
normalising constant survives the tuning. The statement is about $\Zhat$, not $\log\Zhat$, and it holds for every data
set, with or without a true model.

\begin{thm}[The rate cannot be improved; simplified, full version in spec.tex, Theorem~6]\label{thm:sharp}
\textup{(a)} \emph{Every design.} For every $n\ge3$, every standardised design with $p\ge2$, every response $y$ and
every model $S$ with $\abs{S}\le\min(K-1,p/2)$,
\[
  \pis{S}\ \le\ \Big(1+\frac{\abs{S}+1}{3\sqrt n}\Big)^{-1}.
\]

\textup{(b)} \emph{Designs that attain the rate.} Assume A1 and A3. Let the true columns of the design and $\tstar$ be
deterministic, with $\norm{X\tstar}^2\le nM_\mu$ and $\norm{\tstar}^2\le nM_\theta$ for fixed $M_\mu$, $M_\theta$.
Assume one of: \textup{(i)} the null columns of $X$ are orthogonal to each other and to the true columns; or
\textup{(ii)} the raw null columns have independent $\N(0,1)$ entries, independent of the noise, and $p\le n^A$ for a
fixed $A$. Then, for every $p>\sz$, bounded or growing,
\[
  \E\big[1-\pis{\Sz}\big]\ \ge\ (1-o(1))\,(\sz+1)\,\frac{2(p-\sz)}{2p-1-\sz}\,
  \sqrt{\frac{\log\!\big(p\sqrt n/(\sz+1)\big)}{\pi n}}\qquad(\pi=3.14\ldots),
\]
and, if $p\to\infty$, $1-\pis{\Sz}\ge(1-\oP(1))\,(\sz+1)\sqrt{\log p/(\pi n)}$.
\end{thm}
\means Part (a) holds for every data set and has no assumption: under this prior no model, the true one included, can
have posterior mass closer to one than about $(\abs{S}+1)/(3\sqrt n)$. So no design gives a rate faster than
$n^{-1/2}$, and the rate $\rn$ of Theorem~\ref{thm:post} is possible only when $\sz=O(\sqrt{\log n})$; this is why A3
is nearly necessary. Part (b) shows that in the most favourable designs, with orthogonal or independent null
predictors, the posterior still misses the true model by a constant times $\sqrt{\log p/n}$. With $p=n^A$ this is the
rate of Theorem~\ref{thm:post}, logarithm included.

\subsection{Planned}

None of the following is proved. Each is a target, and the choices still open are named.

\begin{planned}[Standard samplers get trapped where ours does not]
There is a family of designs whose posterior has several well-separated modes, on which (i) every Markov chain that
adds or removes one predictor at a time, and the weighted tempered Gibbs sampler of Zanella and Roberts (2019), has a
mixing time that grows quickly with the size of the problem, while (ii) one of Theorems~\ref{thm:post},
\ref{thm:reported} or~\ref{thm:screen} holds on the same designs. Open: the family of designs, which of the three
theorems, the exact class of chains, and the growth rate.
\end{planned}

\begin{planned}[The theory at $p=10^7$]
Take the autoregressive design of the planned $p=10^7$ run at $n=200$, continued as a sequence with $p_n=\lceil
n^{A}\rceil$ and $A=\log(10^7)/\log200\approx3.04$. Show that it satisfies the assumptions of our theorems with this
$A$; for Theorem~\ref{thm:reported} this means exhibiting a leakage set $Q$ that satisfies A9. Open: which theorems
are covered, the set $Q$, and how the sequence is continued. The run itself uses the parallel version of the sampler,
which no theorem covers.
\end{planned}

\begin{planned}[Theorem~\ref{thm:reported} on the raw scale]
Theorem~\ref{thm:reported} takes the coefficients fixed on the standardised scale, and Theorem~\ref{thm:dist} takes
them fixed on the raw scale; with $\sz\ge1$ no design satisfies both. The target is to restate A5(b), A8 and A9 for
fixed raw coefficients and to show, under A7 with $s\ge1$, that with probability tending to one every completed run
returns a list of models $\mathcal T$ with posterior mass $\pi(\mathcal T\mid\Dn)\ge n^{-c_HM/2}$. The Hellinger bound
of Theorem~\ref{thm:dist} would then hold for the posterior restricted to the list the software returns. Open: the
raw-scale form of the assumptions, and whether to aim for the full conclusion $\Shat=\Sz$ on the raw scale.
\end{planned}

\begin{planned}[Theorem~\ref{thm:reported} under an aggregated leak]
The conclusions of Theorem~\ref{thm:reported} hold when the bound A9(a)(i), which is on each predictor outside $Q$
separately, is replaced by a bound on the leaks of all predictors outside $Q$ taken together. Open: the form of the
aggregated condition.
\end{planned}

\begin{planned}[A wrong model, with $p$ fixed]
Drop the linear model altogether. Let the rows $(x_i,y_i)$ be independent draws from any fixed law on
$\R^p\times\R$ with $p$ fixed and $\Sigma=\Var(x)$ positive definite. Let $\beta^0=\Sigma^{-1}\Cov(x,y)$ and
$\alpha^0=\E y-\E x^\top\beta^0$ be the best linear predictor, $r=y-\alpha^0-x^\top\beta^0$ its residual, $v=\E
r^2>0$, and $S^0=\{j:\beta^0_j\neq0\}$. Let $f_0$ be the joint density with $x$ from its true law and $y\mid
x\sim\N(\alpha^0+x^\top\beta^0,v)$. Assume $\E[\norm{x}^2r^2]<\infty$ and $\E r^4<\infty$. Target:
\[
  1-\pi(S^0\mid\Dn)=\OP(n^{-1/2}),\qquad \max_j\big|\PIP_j-\ind{j\in S^0}\big|=\OP(n^{-1/2}),
\]
$\P(\HPM=\MPM=S^0)\to1$, and, for every $M_n\to\infty$,
\[
  \Pi\Big(d_H(f,f_0)>\frac{M_n}{\sqrt n}\ \text{ or }\ \mathrm{KL}(f_0\,\|\,f)>\frac{M_n}{n}\ \Big|\ \Dn\Big)\to0\quad\text{in probability.}
\]
Here $f_0$ is the Kullback--Leibler projection of the true law onto the Gaussian linear family, whenever that
divergence is finite for some member of the family. So without Gaussian errors, a linear conditional mean or constant
error variance, the posterior would select the support of the best linear predictor and concentrate on the closest
Gaussian linear model. The target is for fixed $p$ only.
\end{planned}

\subsection*{References}
{\small
\noindent Song, Q. and Liang, F. (2015). A split-and-merge Bayesian variable selection approach for ultrahigh
dimensional regression. \emph{J. R. Stat. Soc. B} 77, 947--972.

\smallskip
\noindent Zanella, G. and Roberts, G. (2019). Scalable importance tempering and Bayesian variable selection.
\emph{J. R. Stat. Soc. B} 81, 489--517.}

\end{document}
